% lint-skip: 旧jarticle形式のまま未移行 (Issue #593)
\documentclass[a4j]{jarticle}
\usepackage{amsmath}
\usepackage{ascmac}
\usepackage{amssymb}
\usepackage{enumerate}
\usepackage{multicol}
\usepackage{framed}
\usepackage{latexsym}
\usepackage{indent}
\usepackage{cases}
\allowdisplaybreaks
\title{}
\begin{document}
\input{macros}
\begin{oframed}
  $x_1,x_2\dots,x_n$はおのおの$0$，$1$，$2$のどれかの値を取る．$f_1=\dsum{i=1}{n}x_i$，
  $f_2=\dsum{i=1}{n}x_i^2$のとき$f_k=\dsum{i=1}{n}x_i^k(k=1,2,3,\cdots)$を$f_1$と$f_2$とを用いて
  表せ．
\end{oframed}

\setlength{\columnseprule}{0.4pt}
\begin{multicols}{2}
  {\bf[解]}$x_i(i=1,2,\dots,n)$の中に$0$が$a$個，$1$が$b$個，$2$が$c$個あるとすると，
  題意から
  \begin{align}
    a+b+c=n \\
    f_1=b+2c \label{eq:1}\\
    f_2=b+4c  \label{eq:2}\\
    f_k=b+2^kc \label{eq:3}
  \end{align}
  である．\eqref{eq:1}，\eqref{eq:2}から$b=2f_1-f_2$，$c=\dfrac{f_2-f_1}{2}$だから，これを\eqref{eq:3}に代入して
  \begin{align}
    f_k&=2f_1-f_2+2^k\dfrac{f_2-f_1}{2} \\
    &=(2-2^{k-1})f_1+(-1+2^{k-1})f_2\cdots\text{(答)}
  \end{align}
  である．
  \newpage
\end{multicols}
\end{document}
